\section{Training}
\label{sec:training}

The main objective is the usual autoregressive cross-entropy:
\begin{equation}
  \mathcal{L}_{LM}
  = - \sum_t \log p_\theta(x_{t+1}\mid x_{\leq t}).
\end{equation}
Geometric regularization terms are added to prevent collapse of the learned DRM
structure:
\begin{equation}
  \mathcal{L}
  =
  \mathcal{L}_{LM}
  + \lambda_m \mathcal{L}_{metric}
  + \lambda_d \mathcal{L}_{div}
  + \lambda_o \mathcal{L}_{ortho}
  + \lambda_v \mathcal{L}_{var}
  + \lambda_T \mathcal{L}_{torus}.
\end{equation}

The metric regularization keeps $G(x)$ close to the identity early in training.
The diversity loss encourages variation in $U(x)$ across positions. The
orthogonality term prevents rank collapse, while the manifold variance term
prevents all coordinates from contracting to a point.

The toroidal loss is optional and inductive. It groups four manifold dimensions
into two circular pairs and encourages a target radius, angular coverage,
isotropy between cycles, and independence between the two angular variables.
This loss is used as an experimental probe: if the toroidal signature also
appears in weaker or annealed regimes, evidence for spontaneous topological
organization becomes stronger.

Training is implemented in \texttt{scripts/train\_distributed.py}, with support
for single-GPU, DDP, gradient accumulation, mixed precision, checkpointing,
deterministic seeds, and run manifests.
